\documentclass[12pt,a4paper]{amsart}
\usepackage{amsmath,amssymb,amsthm}
\usepackage{mathtools}
\usepackage[utf8]{inputenc}
\usepackage[T1]{fontenc}
\usepackage[ngerman]{babel}
\usepackage{hyperref}
\usepackage{enumitem,booktabs,array}
\usepackage{geometry}
\geometry{margin=2.5cm}

\author{Michael Fuhrmann}
\address{Specialist Mathematics Project, Build 141}

\newtheorem{theorem}{Satz}[section]
\newtheorem{lemma}[theorem]{Lemma}
\newtheorem{korollar}[theorem]{Korollar}
\newtheorem{proposition}[theorem]{Proposition}
\theoremstyle{definition}
\newtheorem{definition}[theorem]{Definition}
\newtheorem{beispiel}[theorem]{Beispiel}
\theoremstyle{remark}
\newtheorem{bemerkung}[theorem]{Bemerkung}
\newtheorem{vermutung}[theorem]{Vermutung}

\title{Zufallsmatrixtheorie: Universalität, Eigenwertstatistiken\\
und das Montgomery-Odlyzko-Gesetz}

\begin{abstract}
Die Zufallsmatrixtheorie (ZMT), durch Wigner in den 1950er-Jahren zur
Modellierung von Energieniveaus schwerer Atomkerne eingeführt, hat sich zu
einem der mächtigsten vereinheitlichenden Rahmenprogramme der modernen
Mathematik und mathematischen Physik entwickelt. Diese Arbeit gibt eine
in sich abgeschlossene Darstellung der drei klassischen Gaußschen
Ensembles -- GOE, GUE und GSE -- samt ihrer grundlegenden Spektralstatistiken.
Wir beweisen das Wigner-Halbkreisgesetz mittels der Momentenmethode,
stellen das Wigner-Surmise für Nächster-Nachbar-Abstände vor und präsentieren
Montgomerys Paarkorrelationssatz für die Nullstellen der Riemann-Zeta-Funktion.
Das zentrale Rätsel -- dass Nullstellen der Zetafunktion sich statistisch
wie GUE-Eigenwerte verhalten (das Montgomery-Odlyzko-Gesetz) -- wird
eingehend diskutiert, einschließlich der Hilbert-Pólya-Vermutung als tiefstmöglicher
Erklärung. Darüber hinaus behandeln wir die Tracy-Widom-Verteilung für
den größten Eigenwert, die Universalitätssätze von Erdős-Yau, das
Marchenko-Pastur-Gesetz für Wishart-Matrizen sowie Anwendungen von der
Kernphysik bis zur Statistik großer Korrelationsmatrizen.
\end{abstract}

\maketitle

\tableofcontents

% ============================================================
\section{Einleitung}
% ============================================================

Die Untersuchung großer Zufallsmatrizen geht auf die bahnbrechenden Arbeiten
von Eugene Wigner \cite{Wigner1955} zurück. Wigner beobachtete, dass die
statistischen Fluktuationen der Energieniveaus schwerer Atomkerne zwar
außerordentlich komplex sind, aber universellen Mustern folgen, die unabhängig
von der spezifischen Kernwechselwirkung sind. Der Schlüsselgedanke war,
den Hamiltonoperator eines schweren Kerns nicht durch eine detaillierte
mikroskopische Rechnung zu modellieren -- was undurchführbar ist -- sondern
durch eine \emph{zufällige} hermitesche Matrix aus einem geeignet invarianten
Wahrscheinlichkeitsensemble. Wigner zeigte, dass im Grenzfall großer Matrizen
die empirische Verteilung der Eigenwerte gegen ein universelles Halbkreisgesetz
konvergiert.

Die Geschichte hätte ein faszinierendes Kapitel der Kernphysik bleiben können,
wäre da nicht eine bemerkenswerte Entdeckung: Die Nullstellen der
Riemann-Zeta-Funktion $\zeta(s)$ scheinen dieselben statistischen Fluktuationseigenschaften
zu besitzen wie die Eigenwerte großer zufälliger unitärer Matrizen (GUE).
Diese als \textbf{Montgomery-Odlyzko-Gesetz} bekannte Beobachtung wurde
auf Grundlage von Montgomerys Paarkorrelationssatz von 1973 \cite{Montgomery1973}
vermutet und von Odlyzko \cite{Odlyzko1987, Odlyzko2001} numerisch mit
außerordentlicher Präzision für zig Millionen Nullstellen bestätigt.

Der Zusammenhang zwischen Spektralstatistiken von Zufallsmatrizen und der
Verteilung der Primzahlen über die Nullstellen von $\zeta(s)$ gehört zu den
tiefsten ungelösten Rätseln der Mathematik. Er legt die Existenz eines
natürlichen Operators -- des \textbf{Hilbert-Pólya-Operators} -- nahe,
dessen Spektrum die Menge der Imaginärteile der nicht-trivialen Nullstellen
von $\zeta(s)$ ist. Die Existenz eines solchen selbstadjungierten Operators
würde unmittelbar die Riemann-Hypothese implizieren.

\subsection*{Überblick über die Arbeit}
Abschnitt~\ref{sec:ensembles} führt die drei Gaußschen Ensembles GOE, GUE
und GSE mit ihren Invarianzeigenschaften ein. Abschnitt~\ref{sec:halbkreis}
beweist das Wigner-Halbkreisgesetz. Abschnitt~\ref{sec:abstaende} behandelt
Nächster-Nachbar-Abstände und das Wigner-Surmise.
Abschnitt~\ref{sec:montgomery} stellt Montgomerys Paarkorrelationssatz vor.
Abschnitt~\ref{sec:goe} behandelt das GOE und die Dyson-Mehta-Identitäten.
Abschnitt~\ref{sec:tracywidom} führt die Tracy-Widom-Verteilung ein.
Abschnitt~\ref{sec:universitaet} gibt einen Überblick über Universalitätssätze.
Abschnitt~\ref{sec:rh} diskutiert den Zusammenhang zur Riemann-Hypothese.
Abschnitt~\ref{sec:anwendungen} behandelt Anwendungen.

% ============================================================
\section{Zufallsmatrix-Ensembles}
\label{sec:ensembles}
% ============================================================

\subsection{Die drei Gaußschen Ensembles}

Die fundamentalen Ensembles der Zufallsmatrixtheorie werden durch die
Symmetrieklasse des Matrixraums und den zugehörigen Dyson-Index $\beta$
unterschieden.

\begin{definition}[Gaußsche Ensembles]
Sei $n \geq 1$ die Matrixdimension.
\begin{enumerate}[label=(\roman*)]
  \item \textbf{Gaußsches Orthogonales Ensemble (GOE)}, $\beta = 1$:
    Das Wahrscheinlichkeitsmaß auf dem Raum der reellen symmetrischen
    $n \times n$-Matrizen $H = H^T$ mit Dichte proportional zu
    \[
      P(H)\,dH \;\propto\; \exp\!\bigl(-\tfrac{n}{4}\,\operatorname{Sp}(H^2)\bigr)\,dH,
    \]
    wobei $dH = \prod_{i \leq j} dH_{ij}$.

  \item \textbf{Gaußsches Unitäres Ensemble (GUE)}, $\beta = 2$:
    Das Wahrscheinlichkeitsmaß auf dem Raum der $n \times n$-hermiteschen
    Matrizen $H = H^\dagger$ mit Dichte proportional zu
    \[
      P(H)\,dH \;\propto\; \exp\!\bigl(-\tfrac{n}{2}\,\operatorname{Sp}(H^2)\bigr)\,dH.
    \]

  \item \textbf{Gaußsches Symplektisches Ensemble (GSE)}, $\beta = 4$:
    Das Wahrscheinlichkeitsmaß auf dem Raum der $n \times n$-selbstdualen
    quaternionischen Matrizen $H = H^R$ mit Dichte proportional zu
    \[
      P(H)\,dH \;\propto\; \exp\!\bigl(-n\,\operatorname{Sp}(H^2)\bigr)\,dH.
    \]
\end{enumerate}
\end{definition}

\subsection{Invarianzeigenschaften und physikalische Interpretation}

Eine entscheidende Eigenschaft dieser Ensembles ist ihre Invarianz:
Das GOE ist invariant unter orthogonaler Konjugation $H \mapsto O H O^T$
($O \in O(n)$), das GUE unter unitärer Konjugation $H \mapsto U H U^\dagger$
($U \in U(n)$), und das GSE unter symplektischer Konjugation.

Physikalisch kodiert der Dyson-Index $\beta$ die Symmetrie des zugrundeliegenden
Quantensystems:
\begin{itemize}
  \item $\beta = 1$ (GOE): Systeme mit \emph{Zeitumkehrinvarianz} und
    ganzzahligem Spin (z.\,B.\ schwere Kerne, Quantenbillards ohne Magnetfeld).
  \item $\beta = 2$ (GUE): Systeme \emph{ohne} Zeitumkehrinvarianz
    (z.\,B.\ Elektronen im Magnetfeld, komplexes Quantenchaos).
  \item $\beta = 4$ (GSE): Systeme mit Zeitumkehrinvarianz und halbzahligem Spin
    (Kramers-Entartung).
\end{itemize}

\subsection{Gemeinsame Eigenwertdichte}

Ein Eckpfeilerresultat ist die explizite Form der gemeinsamen
Wahrscheinlichkeitsdichte der Eigenwerte nach Integration über die
Eigenvektorfreiheitsgrade.

\begin{theorem}[Gemeinsame Eigenwertdichte]
\label{thm:jpdf}
Für jedes der drei Gaußschen Ensembles mit Parameter $\beta \in \{1,2,4\}$
ist die gemeinsame Wahrscheinlichkeitsdichte der Eigenwerte
$\lambda_1 \leq \cdots \leq \lambda_n$ gegeben durch
\[
  P_\beta(\lambda_1,\ldots,\lambda_n)
  \;=\; C_{n,\beta}\;
  \prod_{1 \leq i < j \leq n} |\lambda_i - \lambda_j|^\beta\;
  \exp\!\Bigl(-\frac{n\beta}{4}\sum_{i=1}^n \lambda_i^2\Bigr),
\]
wobei $C_{n,\beta}$ eine explizite Normierungskonstante ist, die durch das
Selberg-Integral gegeben wird.
\end{theorem}

Der Vandermonde-artige Faktor $|\lambda_i - \lambda_j|^\beta$ ist entscheidend:
Er erzeugt \textbf{Niveau-Abstoßung} -- Eigenwerte meiden einander, wobei
die Stärke der Abstoßung mit $\beta$ wächst.

% ============================================================
\section{Wigners Halbkreisgesetz}
\label{sec:halbkreis}
% ============================================================

\subsection{Formulierung des Satzes}

Das Halbkreisgesetz ist das Analogon des Zentralen Grenzwertsatzes für
das empirische Spektralmaß großer Zufallsmatrizen.

\begin{theorem}[Wigners Halbkreisgesetz, 1958]
\label{thm:halbkreis}
Sei $H_n$ eine $n \times n$-hermitesche Zufallsmatrix, deren Einträge
$(H_{ij})_{i \leq j}$ unabhängig (bis auf die hermitesche Bedingung) sind
mit Erwartungswert null und Varianz $\sigma^2 < \infty$. Betrachte die
umskalierte Matrix $\tilde{H}_n = H_n / (\sigma\sqrt{n})$. Dann konvergiert
das empirische Spektralmaß
\[
  \mu_n \;=\; \frac{1}{n}\sum_{k=1}^n \delta_{\lambda_k(\tilde{H}_n)}
\]
schwach in Wahrscheinlichkeit gegen die Wigner-Halbkreisverteilung
\[
  \rho_{\mathrm{HK}}(x) \;=\; \frac{1}{2\pi}\sqrt{(4-x^2)_+},
\]
getragen auf $[-2,2]$. Dabei ist $(y)_+ = \max(y,0)$.
In der auf das Einheitsintervall normierten Form:
\[
  \rho(x) \;=\; \frac{2}{\pi}\sqrt{1-x^2}, \quad |x| \leq 1.
\]
\end{theorem}

\begin{bemerkung}
Die Formel $\rho_{\mathrm{HK}}(x) = \frac{1}{2\pi}\sqrt{4-x^2}$ (getragen
auf $[-2,2]$) entsteht natürlich aus der im GUE verwendeten Normierung
$\operatorname{Var}[H_{ij}] = 1/n$ mittels $H = (A + A^\dagger)/\sqrt{2n}$.
Die beiden Konventionen sind durch die Substitution $x \mapsto 2x$ verbunden.
\end{bemerkung}

\subsection{Beweis mittels Momentenmethode}

Wir skizzieren den Beweis mit der Momentenmethode, dem elementarsten Ansatz,
der zugleich die exakte Universalitätsaussage liefert.

\begin{proof}[Beweissskizze]
Das $k$-te Moment von $\mu_n$ ist
\[
  m_k^{(n)} \;=\; \int x^k\,d\mu_n(x)
  \;=\; \frac{1}{n}\operatorname{Sp}(\tilde{H}_n^k).
\]
Wir berechnen $\mathbb{E}[m_k^{(n)}]$ durch Entwicklung der Spur:
\[
  \mathbb{E}\!\left[\frac{1}{n}\operatorname{Sp}(H_n^k)\right]
  = \frac{1}{n \cdot n^{k/2}} \sum_{i_1,\ldots,i_k=1}^n
    \mathbb{E}[H_{i_1 i_2} H_{i_2 i_3} \cdots H_{i_k i_1}].
\]
Da die Einträge Erwartungswert null haben, verschwindet der Erwartungswert
eines Produkts, sofern nicht jedes Indexpaar $(i_j, i_{j+1})$ mindestens
zweimal auftritt. In führender Ordnung $n \to \infty$ überleben nur
\emph{nicht-kreuzende Paarungen} (auch \emph{Catalan-Paarungen} genannt),
und jede trägt $\sigma^{2k/2} \cdot n^{k/2+1}$ bei.
Die Anzahl solcher Paarungen für $k = 2m$ ist die Catalan-Zahl
\[
  C_m \;=\; \frac{1}{m+1}\binom{2m}{m}.
\]
Diese sind genau die geraden Momente der Halbkreisverteilung:
$\int x^{2m}\rho_{\mathrm{HK}}(x)\,dx = C_m$.
Die ungeraden Momente verschwinden aus Symmetriegründen. Da die Momente
die Halbkreisverteilung charakterisieren (Carlemans Bedingung ist erfüllt),
folgt die Konvergenz $\mu_n \to \rho_{\mathrm{HK}}$.
\end{proof}

\subsection{Universalität des Halbkreisgesetzes}

Der obige Beweis zeigt, dass das Halbkreisgesetz weit über den Gaußschen
Fall hinaus gilt: Es gilt für beliebige Wigner-Matrizen, deren Einträge
Erwartungswert null und endliche Varianz haben. Die Gaußsche Struktur ist
für die globale Dichte nicht erforderlich.

% ============================================================
\section{Eigenwertabstandsstatistiken}
\label{sec:abstaende}
% ============================================================

\subsection{Nächster-Nachbar-Abstände und Niveau-Abstoßung}

Während das Halbkreisgesetz die \emph{globale} Eigenwertdichte beschreibt,
ist die Zufallsmatrixtheorie am bekanntesten für ihre Vorhersagen über
\emph{lokale} Statistiken -- insbesondere die Verteilung der Abstände
zwischen benachbarten Eigenwerten.

Seien $\lambda_1 \leq \lambda_2 \leq \cdots \leq \lambda_n$ die Eigenwerte
einer Zufallsmatrix aus dem GUE (oder GOE). Man bildet die normierten
Nächster-Nachbar-Abstände
\[
  s_k \;=\; \frac{\lambda_{k+1} - \lambda_k}{\langle \lambda_{k+1} - \lambda_k \rangle},
\]
normiert auf mittleren Abstand 1. Die empirische Verteilung der $s_k$
konvergiert für $n \to \infty$ gegen eine Grenzverteilung.

\subsection{Das Wigner-Surmise}

Für $2 \times 2$-Matrizen berechnete Wigner \cite{Wigner1957} die exakte
Abstandsverteilung, die sich als hervorragende Näherung für $n \times n$-Matrizen
mit großem $n$ herausstellt. Dies ist das \textbf{Wigner-Surmise}.

\begin{theorem}[Wigner-Surmise]
\label{thm:surmise}
Die Nächster-Nachbar-Abstandsverteilung der Eigenwerte des $n \times n$-GUE
konvergiert nach Entfaltung (Normierung auf mittleren Abstand 1) gegen das
Wigner-Surmise für $\beta = 2$:
\[
  p_{\mathrm{GUE}}(s) \;=\; \frac{32}{\pi^2}\,s^2\,\exp\!\left(-\frac{4s^2}{\pi}\right),
  \quad s \geq 0.
\]
Für das GOE ($\beta = 1$):
\[
  p_{\mathrm{GOE}}(s) \;=\; \frac{\pi}{2}\,s\,\exp\!\left(-\frac{\pi s^2}{4}\right),
  \quad s \geq 0.
\]
Für das GSE ($\beta = 4$):
\[
  p_{\mathrm{GSE}}(s) \;=\; \frac{2^{18}}{3^6\pi^3}\,s^4\,\exp\!\left(-\frac{64 s^2}{9\pi}\right),
  \quad s \geq 0.
\]
\end{theorem}

\subsection{Niveau-Abstoßung}

Ein fundamentales Merkmal aller drei Formeln ist $p_\beta(0) = 0$:
Eigenwerte \textbf{stoßen sich ab}. Das Verhalten nahe null ist
\[
  p_\beta(s) \;\sim\; c_\beta\, s^\beta, \quad s \to 0^+,
\]
die Abstoßung ist also \emph{linear} für GOE, \emph{quadratisch} für GUE
und \emph{quartisch} für GSE. Dies steht im scharfen Kontrast zur
Poisson-Verteilung $p_{\mathrm{Poisson}}(s) = e^{-s}$, für die $p(0) = 1$,
was \emph{keiner Abstoßung} entspricht (Eigenwerte sind unabhängig).

Der physikalische Mechanismus: Beim GOE ($\beta=1$) erfordert eine Matrix
mit entartetem Eigenwert eine Bedingung der Kodimension 1 (eine reelle
Gleichung), was Entartungen selten macht. Beim GUE ($\beta=2$) sind es zwei
Bedingungen (eine komplexe Gleichung), daher eine höhere Potenz von $s$.

% ============================================================
\section{Montgomerys Paarkorrelation}
\label{sec:montgomery}
% ============================================================

\subsection{Die Paarkorrelation der Zeta-Nullstellen}

Sei $\rho_n = \frac{1}{2} + i\gamma_n$ die nicht-triviale Nullstellen von
$\zeta(s)$, wobei wir für diesen Abschnitt die Riemann-Hypothese (RH)
annehmen, sodass alle $\gamma_n$ reell und positiv sind. Die Nullstellen
sind geordnet: $0 < \gamma_1 \leq \gamma_2 \leq \cdots$ mit
$\gamma_1 \approx 14{,}135$.

In Analogie zu Zufallsmatrix-Eigenwerten entfaltet man die Nullstellenfolge
zunächst durch Anwendung der lokalen mittleren Dichte. Aus der expliziten
Formel und dem Primzahlsatz folgt
\[
  \gamma_n \;\approx\; \frac{2\pi n}{\log(n/2\pi e)}.
\]
Nach Entfaltung (Setzen $\tilde{\gamma}_n = \frac{\gamma_n}{2\pi}\log\frac{\gamma_n}{2\pi}$)
haben die normierten Abstände Mittelwert 1.

\begin{theorem}[Montgomerys Paarkorrelationssatz, 1973]
\label{thm:montgomery}
Unter Annahme der Riemann-Hypothese gilt für jede gerade Schwartz-Funktion $f$,
deren Fouriertransformierte $\hat{f}$ in $(-1,1)$ getragen ist:
\[
  \frac{1}{N(T)}\sum_{\substack{m,n \\ m \neq n}} f\!\left(\frac{\gamma_m - \gamma_n}{2\pi/\log T}\right)
  \;\longrightarrow\;
  \int_{-\infty}^\infty f(r)\,R_2(r)\,dr \quad \text{für } T \to \infty,
\]
wobei die Paarkorrelationsfunktion
\[
  R_2(r) \;=\; 1 - \left(\frac{\sin(\pi r)}{\pi r}\right)^2
\]
ist. Hierbei ist $N(T) = \#\{\gamma_n : 0 < \gamma_n \leq T\}$ die Nullstellen-Zählfunktion.
\end{theorem}

\begin{bemerkung}
Die Funktion $R_2(r) = 1 - (\sin(\pi r)/\pi r)^2$ ist \emph{exakt} die
Zweipunkt-Korrelationsfunktion des GUE (Eigenwerte großer unitärer Zufallsmatrizen).
Montgomerys Satz beweist damit, dass die Paarkorrelation der $\zeta$-Nullstellen
mit der der GUE-Eigenwerte übereinstimmt,
\emph{im eingeschränkten Bereich, wo $\hat{f}$ in $(-1,1)$ getragen ist}.
Die Erweiterung auf alle Testfunktionen und die vollen GUE-Statistiken
(einschließlich aller $n$-Punkt-Korrelationen) für $\zeta$-Nullstellen bleiben
eine offene Vermutung.
\end{bemerkung}

\subsection{Odlyzkos numerische Verifikation}

Andrew Odlyzko \cite{Odlyzko1987} berechnete die ersten $10^6$ Nullstellen
nahe $\gamma \approx 10^{20}$ mit hoher Genauigkeit und erstellte Histogramme
der normierten Abstände. Die Übereinstimmung mit dem GUE-Wigner-Surmise
$p_{\mathrm{GUE}}(s)$ ist bemerkenswert. Die Poisson-Verteilung ist klar
ausgeschlossen.

In einer späteren Arbeit \cite{Odlyzko2001} berechnete Odlyzko Nullstellen nahe Höhe $10^{22}$
und bestätigte die GUE-Statistiken mit noch höherer Genauigkeit.
Diese Berechnungen stellen den stärksten numerischen Beleg für die Riemann-Hypothese dar.

% ============================================================
\section{Gaußsches Orthogonales Ensemble}
\label{sec:goe}
% ============================================================

\subsection{GOE und Zeitumkehrinvarianz}

Das GOE ($\beta = 1$) ist das geeignete Ensemble für Quantensysteme mit
Zeitumkehrinvarianz und ganzzahligem (oder null) Spin. Das kanonische
Beispiel ist ein Billardsystem mit klassisch chaotischer Geometrie
(Sinai-Billard, Bunimovich-Stadion).

Für das GOE haben die Eigenwerte $\lambda_1 \leq \cdots \leq \lambda_n$
die gemeinsame Dichte (Satz~\ref{thm:jpdf} mit $\beta = 1$):
\[
  P_1(\lambda_1,\ldots,\lambda_n)
  \;=\; C_{n,1}\;\prod_{i<j}|\lambda_i - \lambda_j|\;\exp\!\Bigl(-\frac{n}{4}\sum_i \lambda_i^2\Bigr).
\]

\subsection{Niveauzahlvarianz und die Dyson-Mehta-Formel}

Eine wichtige Statistik ist die \textbf{Niveauzahlvarianz} $\Sigma^2(L)$,
definiert als die Varianz der Anzahl von Eigenwerten in einem Fenster der
Länge $L$ (in entfalteten Koordinaten). Für das GOE bewiesen Dyson und
Mehta \cite{Mehta2004}:

\begin{theorem}[Dyson-Mehta-Identität für das GOE]
\label{thm:dysonmehta}
Für das GOE im Grenzfall $n \to \infty$ gilt
\[
  \Sigma^2(L) \;=\; \frac{2}{\pi^2}\left(\log(2\pi L) + \gamma_{\mathrm{E}} + 1 - \frac{\pi^2}{8}\right)
  + O\!\left(\frac{1}{L}\right),
\]
wobei $\gamma_{\mathrm{E}}$ die Euler-Mascheroni-Konstante ist. Im Gegensatz dazu
gilt für Poisson: $\Sigma^2_{\mathrm{Poisson}}(L) = L$.
Das logarithmische Wachstum von $\Sigma^2$ (statt linear für Poisson) spiegelt
die starke \emph{langreichweitige Starrheit} des GOE-Eigenwertspektrums wider.
\end{theorem}

\subsection{GOE-Niveauabstand: die lineare Abstoßung}

Wie im Wigner-Surmise (Satz~\ref{thm:surmise}) formuliert, zeigt die GOE-Abstandsverteilung
$p_{\mathrm{GOE}}(s) = (\pi/2)\,s\,e^{-\pi s^2/4}$ lineare Niveau-Abstoßung:
$p_{\mathrm{GOE}}(s) \sim (\pi/2)\,s$ für $s \to 0$.
Dies ist charakteristisch für Quantensysteme mit Zeitumkehrinvarianz und
wurde in Hunderten von Experimenten mit Kern-, Atom- und Molekülspektren bestätigt.

% ============================================================
\section{Tracy-Widom-Verteilung}
\label{sec:tracywidom}
% ============================================================

\subsection{Der größte Eigenwert}

Die \textbf{Tracy-Widom-Verteilungen} beschreiben die Fluktuationen des
größten Eigenwerts $\lambda_{\max}$ einer Zufallsmatrix um seine erwartete
Position am Rand des Spektrums.

Für das GUE mit $n \times n$-Matrizen konzentriert sich $\lambda_{\max}$
um den rechten Rand des Halbkreises bei $2$. Der Tracy-Widom-Satz liefert
die präzise Grenzverteilung der Fluktuationen auf der Skala $n^{-2/3}$.

\begin{theorem}[Tracy-Widom, 1994]
\label{thm:tracywidom}
Sei $H_n$ aus dem GUE mit Standardnormierung
$\mathbb{E}[|H_{ij}|^2] = 1/n$ entnommen. Dann gilt
\[
  \mathbb{P}\!\left(\frac{\lambda_{\max} - 2}{n^{-2/3}} \leq s\right)
  \;\longrightarrow\; F_2(s) \quad \text{für } n \to \infty,
\]
wobei $F_2$ die \textbf{GUE-Tracy-Widom-Verteilung} ist:
\[
  F_2(s) \;=\; \exp\!\left(-\int_s^\infty (x-s)\,q(x)^2\,dx\right).
\]
Dabei ist $q(x)$ die eindeutige Lösung der \textbf{Hastings-McLeod-Gleichung}
\[
  q'' \;=\; xq + 2q^3, \quad q(x) \sim \mathrm{Ai}(x) \text{ für } x \to +\infty,
\]
wobei $\mathrm{Ai}$ die Airy-Funktion ist.
\end{theorem}

\begin{bemerkung}
Analoge Resultate gelten für GOE und GSE:
\begin{align*}
  F_1(s) &= \exp\!\left(-\frac{1}{2}\int_s^\infty q(x)\,dx\right)\cdot [F_2(s)]^{1/2}
  \quad \text{(GOE-Tracy-Widom)},\\
  F_4(s/\sqrt{2}) &= \cosh\!\left(\frac{1}{2}\int_s^\infty q(x)\,dx\right)\cdot [F_2(s)]^{1/2}
  \quad \text{(GSE-Tracy-Widom)}.
\end{align*}
\end{bemerkung}

\subsection{Anwendungen der Tracy-Widom-Verteilung}

Die Tracy-Widom-Verteilungen erscheinen weit über die Zufallsmatrixtheorie hinaus:
\begin{itemize}
  \item Die Länge der längsten aufsteigenden Teilfolge einer zufälligen
    Permutation von $\{1,\ldots,n\}$ (Satz von Baik-Deift-Johansson \cite{BaikDeiftJohansson1999}).
  \item Die Höhenfluktuationen zufälliger Wachstumsinterfaces in der
    Kardar-Parisi-Zhang (KPZ)-Universalitätsklasse.
  \item Der größte Singulärwert bei der Hauptkomponentenanalyse (PCA)
    von Rauschmatrizen.
  \item Die freie Energie gerichteter Polymere in zufälligen Medien.
\end{itemize}

% ============================================================
\section{Universalität}
\label{sec:universitaet}
% ============================================================

\subsection{Was Universalität bedeutet}

Eines der zentralen Themen der modernen Zufallsmatrixtheorie ist
\textbf{Universalität}: Die lokalen Eigenwertstatistiken (Niveauabstände,
$n$-Punkt-Korrelationen) hängen nur von der Symmetrieklasse ($\beta = 1, 2, 4$)
des Ensembles ab, nicht von der spezifischen Verteilung der Matrixeinträge.

Dies ist analog zur Universalität der Gaußschen Verteilung im Zentralen
Grenzwertsatz: Die globale Spektraldichte ist universell (der Halbkreis),
und ebenso sind es die lokalen Statistiken (GUE/GOE-Korrelationen).

\subsection{Der Erdős-Yau-Universalitätssatz}

Das folgende Durchbruchsergebnis erweiterte die GUE/GOE-Universalität weit
über den Gaußschen Fall hinaus:

\begin{theorem}[Erdős-Schlein-Yau-Yin, 2010--2012]
\label{thm:universitaet}
Sei $H_n$ eine Wigner-Matrix (d.\,h.\ hermitesch mit unabhängigen oberen
Dreieckseinträgen mit Erwartungswert null und Varianz $1/n$), die eine
sub-exponentielle Endbedingung erfüllt. Dann gilt:
\begin{enumerate}[label=(\roman*)]
  \item Das empirische Spektralmaß konvergiert gegen das Wigner-Halbkreisgesetz
    (lokales Gesetz, Bulk und Rand).
  \item Alle endlichdimensionalen Korrelationsfunktionen der entfalteten Eigenwerte
    im Bulk konvergieren gegen die des GUE (für hermitesche Wigner-Matrizen)
    bzw.\ GOE (für reell-symmetrische Wigner-Matrizen).
  \item Das Tracy-Widom-Gesetz gilt für den größten Eigenwert.
\end{enumerate}
\end{theorem}

\begin{bemerkung}
Der Beweis von Satz~\ref{thm:universitaet} verwendet zwei Hauptwerkzeuge:
(a) ein \emph{lokales Halbkreisgesetz} (Starrheit der Eigenwerte auf der Skala $1/n$)
und (b) das \emph{Dyson-Brownsche-Bewegungs}-Interpolationsargument, das zeigt,
dass die GUE-Statistiken unter kleinen Störungen der Eintragsverteilung
lokal stabil sind.
\end{bemerkung}

\subsection{Beta-Ensembles und log-Gas}

Für allgemeines $\beta > 0$ kann man $\beta$-Ensembles durch die gemeinsame
Eigenwertdichte
\[
  P_\beta(\lambda_1,\ldots,\lambda_n) \;\propto\;
  \prod_{i<j} |\lambda_i - \lambda_j|^\beta \cdot e^{-n\beta V(\lambda_i)/2}
\]
mit einem allgemeinen einschließenden Potential $V$ definieren. Das
$\beta$-Ensemble interpoliert zwischen Poisson ($\beta \to 0$), GOE ($\beta=1$),
GUE ($\beta=2$), GSE ($\beta=4$) und dem ``starren'' Kristall ($\beta \to \infty$).

% ============================================================
\section{Verbindung zur Riemann-Hypothese}
\label{sec:rh}
% ============================================================

\subsection{Die Hilbert-Pólya-Vermutung}

Die tiefstmögliche Erklärung der GUE-Statistiken der $\zeta$-Nullstellen
ist die folgende:

\begin{vermutung}[Hilbert-Pólya-Vermutung]
\label{verm:hilbertpolya}
Es gibt einen selbstadjungierten unbeschränkten Operator $\mathcal{L}$ auf einem
Hilbert-Raum $\mathcal{H}$, so dass das Spektrum von $\mathcal{L}$ genau
\[
  \sigma(\mathcal{L}) \;=\; \left\{ \gamma_n \;:\; \zeta\!\left(\tfrac{1}{2} + i\gamma_n\right) = 0,\; \gamma_n \in \mathbb{R} \right\}
\]
ist.
\end{vermutung}

\noindent
Würde ein solcher Operator existieren und hermitesch (selbstadjungiert) sein,
wären seine Eigenwerte notwendigerweise reell -- was genau der Inhalt der
Riemann-Hypothese ist. Wäre $\mathcal{L}$ zudem von der GUE-Universalitätsklasse
(z.\,B.\ ein Quanten-Hamiltonoperator ohne Zeitumkehrinvarianz), erkläre dies
die GUE-Statistiken.

\begin{bemerkung}
Die Hilbert-Pólya-Vermutung (Vermutung~\ref{verm:hilbertpolya}) ist
\textbf{vollständig offen}. Kein konkreter Kandidat für $\mathcal{L}$ ist
rigoros etabliert, obwohl viele Vorschläge existieren (Schrödinger-Operatoren
auf modularen Flächen, $p$-adische Quantenmechanik, der Berry-Keating-Operator
$xp + px$, usw.). Die Vermutung ist älter als die Zufallsmatrixtheorie und
wurde von Hilbert und Pólya unabhängig um 1910--1920 formuliert.
\end{bemerkung}

\subsection{Das Montgomery-Odlyzko-Gesetz als numerischer Beleg}

\begin{vermutung}[Montgomery-Odlyzko-Gesetz, vollständige Version]
\label{verm:moo}
Alle $n$-Punkt-Korrelationsfunktionen der normierten $\zeta$-Nullstellen
konvergieren gegen die des GUE, nicht nur die Paarkorrelation (die unter RH
für Testfunktionen mit eingeschränktem Fourier-Träger durch
Satz~\ref{thm:montgomery} bewiesen ist).
\end{vermutung}

Der numerische Beleg aus Odlyzkos Berechnungen \cite{Odlyzko1987, Odlyzko2001}
ist außerordentlich überzeugend: Das Wigner-Surmise $p_{\mathrm{GUE}}(s)$
passt das Histogramm der $\zeta$-Abstände über $10^{22}$ Nullstellen
im Rahmen statistischer Fluktuationen.

\subsection{Weitere Evidenz und verwandte Resultate}

\begin{itemize}
  \item \textbf{$L$-Funktionen}: Ähnliche GUE- oder GOE-Statistiken gelten
    für die Nullstellen von Dirichlet-$L$-Funktionen, wobei das Ensemble vom
    Symmetrietyp der $L$-Funktion abhängt (Katz-Sarnak \cite{KatzSarnak1999}).
  \item \textbf{Momente von $\zeta$}: Keating-Snaith \cite{KeatingSnaith2000}
    verwendeten GUE-Charakteristikpolynommomente, um die genaue führende Konstante
    in der Asymptotik $\int_0^T |\zeta(\frac{1}{2}+it)|^{2k}\,dt \sim c_k\,T(\log T)^{k^2}$
    zu vermuten.
  \item \textbf{Quantenchaos}: Die Bohigas-Giannoni-Schmit-Vermutung besagt,
    dass die Energieniveaus eines Quantensystems, dessen klassische Dynamik chaotisch ist,
    GUE- (oder GOE-)Statistiken folgen. Dies bleibt eine Vermutung, gestützt
    durch überwältigende numerische Evidenz.
\end{itemize}

% ============================================================
\section{Anwendungen}
\label{sec:anwendungen}
% ============================================================

\subsection{Kernphysik}

Wigners ursprüngliche Motivation war die Statistik der Energieniveaus schwerer
Kerne ($A > 100$). Der Hamiltonoperator eines solchen Kerns enthält
$\sim 10^{50}$ Matrixelemente, was direkte Berechnungen unmöglich macht.
Das Zufallsmatrixmodell macht präzise statistische Vorhersagen, die in
Hunderten von gemessenen Kernspektren bestätigt wurden, beginnend mit der
Grundlagenanalyse von Bohigas-Haq-Pandey des ``Nuclear Data Ensemble'' (1983).

\subsection{Quantenchaos}

Die Zufallsmatrixtheorie ist das zentrale Werkzeug im Forschungsfeld
\textbf{Quantenchaos}, das untersucht, wie sich klassisches Chaos in der
Quantenmechanik manifestiert. Die zentrale empirische Tatsache
(BGS-Vermutung) lautet:
\begin{quote}
  Quantisierungen klassisch chaotischer Hamiltonoperatoren haben GUE- oder GOE-Niveaustatistiken;
  Quantisierungen klassisch integrabler Hamiltonoperatoren haben Poisson-Statistiken.
\end{quote}
Paradigmatische Beispiele sind das Sinai-Billard (GOE), das Bunimovich-Stadion (GOE)
und Magnetbillards (GUE).

\subsection{Große zufällige Korrelationsmatrizen und Finanzwissenschaft}

In der hochdimensionalen Statistik bildet man häufig eine Stichprobenkovarianzmatrix
$\hat{\Sigma} = \frac{1}{m} X X^T$, wobei $X$ eine $n \times m$-Datenmatrix ist
mit $n, m \to \infty$, $n/m \to \gamma$. Das Marchenko-Pastur-Gesetz
\cite{MarchenkoPastur1967} beschreibt die Grenzspektraldichte:

\begin{theorem}[Marchenko-Pastur-Gesetz]
\label{thm:marchenkopastur}
Sei $X$ eine $n \times m$-Matrix mit unabhängig identisch verteilten Einträgen
mit Erwartungswert null und Varianz $\sigma^2$, und sei $\gamma = n/m$. Dann
konvergiert das empirische Spektralmaß von $\frac{1}{m} X X^T$ gegen die
Marchenko-Pastur-Verteilung
\[
  \rho_{\gamma}(x)
  \;=\; \frac{\sqrt{(\lambda_+ - x)(x - \lambda_-)}}{2\pi\,\gamma\,\sigma^2\,x},
  \quad x \in [\lambda_-, \lambda_+],
\]
mit $\lambda_{\pm} = \sigma^2(1 \pm \sqrt{\gamma})^2$.
\end{theorem}

In der Finanzwirtschaft und den Datenwissenschaften wird dieses Gesetz verwendet,
um in empirischen Korrelationsmatrizen von Aktienrenditen \emph{Signal}
(Eigenwerte oberhalb von $\lambda_+$) von \emph{Rauschen} (Eigenwerte innerhalb
von $[\lambda_-, \lambda_+]$) zu unterscheiden.

\subsection{Drahtlose Kommunikation}

In MIMO (Multiple-Input Multiple-Output)-Funksystemen wird die Kanalkapazität
durch die Singulärwerte der Kanalmatrix bestimmt, die Marchenko-Pastur-Statistiken
folgt, wenn der Kanal zufällig (Rayleigh-Fading) ist. Die Zufallsmatrixtheorie
liefert präzise Vorhersagen für die ergodische Kapazität und
Ausfallwahrscheinlichkeit von MIMO-Systemen.

\subsection{Statistische Mechanik und Wachstumsmodelle}

Die Tracy-Widom-Verteilung und verwandte Resultate verbinden die
Zufallsmatrixtheorie mit:
\begin{itemize}
  \item Der KPZ-Universalitätsklasse zufälliger Wachstumsprozesse (Eden-Modell,
    ASEP, polynukleares Wachstumsmodell).
  \item Freier Wahrscheinlichkeitstheorie (Voiculescu), die klassische
    Wahrscheinlichkeitstheorie auf nicht-kommutative Zufallsvariablen erweitert.
  \item Integrablen Systemen und der Theorie isomonodromischer Deformationen.
\end{itemize}

% ============================================================
\section{Literatur}
\label{sec:literatur}
% ============================================================

\begin{thebibliography}{99}

\bibitem{Wigner1955}
Wigner, E.P.\ (1955).
\textit{Characteristic vectors of bordered matrices with infinite dimensions}.
Annals of Mathematics, 62(3), 548--564.

\bibitem{Wigner1957}
Wigner, E.P.\ (1957).
\textit{Distribution of neutron resonance level spacings}.
Conference on Neutron Physics by Time-of-Flight, Gatlinburg, Tennessee,
ORNL-2309.

\bibitem{Montgomery1973}
Montgomery, H.L.\ (1973).
\textit{The pair correlation of zeros of the zeta function}.
Analytic Number Theory (Proc.\ Sympos.\ Pure Math.\ Vol.\ XXIV),
American Mathematical Society, Providence, RI, S.\ 181--193.

\bibitem{Odlyzko1987}
Odlyzko, A.M.\ (1987).
\textit{On the distribution of spacings between zeros of the zeta function}.
Mathematics of Computation, 48(177), 273--308.

\bibitem{Odlyzko2001}
Odlyzko, A.M.\ (2001).
\textit{The $10^{22}$-nd zero of the Riemann zeta function}.
Dynamical, Spectral, and Arithmetic Zeta Functions (San Antonio, TX, 1999),
Contemporary Mathematics, 290, S.\ 139--144, AMS.

\bibitem{TracyWidom1994}
Tracy, C.A., Widom, H.\ (1994).
\textit{Level-spacing distributions and the Airy kernel}.
Communications in Mathematical Physics, 159(1), 151--174.

\bibitem{Mehta2004}
Mehta, M.L.\ (2004).
\textit{Random Matrices}, 3.\ Auflage.
Pure and Applied Mathematics, Bd.\ 142, Academic Press, Amsterdam.

\bibitem{AndersonGuionnetZeitouni2010}
Anderson, G.W., Guionnet, A., Zeitouni, O.\ (2010).
\textit{An Introduction to Random Matrices}.
Cambridge Studies in Advanced Mathematics, Bd.\ 118,
Cambridge University Press.

\bibitem{ErdosYau2012}
Erdős, L., Yau, H.-T.\ (2012).
\textit{Universality of local spectral statistics of random matrices}.
Bulletin of the American Mathematical Society, 49(3), 377--414.

\bibitem{MarchenkoPastur1967}
Mar\v{c}enko, V.A., Pastur, L.A.\ (1967).
\textit{Distribution of eigenvalues for some sets of random matrices}.
Mathematics of the USSR-Sbornik, 1(4), 457--483.

\bibitem{KatzSarnak1999}
Katz, N.M., Sarnak, P.\ (1999).
\textit{Random Matrices, Frobenius Eigenvalues, and Monodromy}.
American Mathematical Society Colloquium Publications, Bd.\ 45, AMS.

\bibitem{KeatingSnaith2000}
Keating, J.P., Snaith, N.C.\ (2000).
\textit{Random matrix theory and $\zeta(1/2+it)$}.
Communications in Mathematical Physics, 214(1), 57--89.

\bibitem{BaikDeiftJohansson1999}
Baik, J., Deift, P., Johansson, K.\ (1999).
\textit{On the distribution of the length of the longest increasing subsequence
of random permutations}.
Journal of the American Mathematical Society, 12(4), 1119--1178.

\end{thebibliography}

\end{document}
